\section{ОБЗОР ПРЕДМЕТНОЙ ОБЛАСТИ}

\subsection{Мульти-роботизированные системы}

Мульти-роботизированная система (МРС) представляет собой совокупность нескольких автономных роботов, функционирующих в общей среде и взаимодействующих между собой для достижения заданных целей~\cite{parker2008mrs}. В отличие от одиночных роботизированных решений, МРС обладают рядом существенных преимуществ.

Во-первых, \textit{параллелизм}: несколько роботов способны одновременно выполнять различные задачи, что сокращает общее время выполнения миссии. Во-вторых, \textit{отказоустойчивость}: выход из строя одного агента не приводит к полной остановке системы, поскольку оставшиеся роботы могут перераспределить задачи. В-третьих, \textit{масштабируемость}: производительность системы может быть увеличена добавлением новых роботов без изменения управляющей архитектуры.

Практические применения МРС охватывают широкий спектр областей. В складской логистике системы автономных мобильных роботов обеспечивают перемещение товаров между стеллажами и зонами комплектации. Компания Amazon (ранее Kiva Systems) продемонстрировала возможность координации сотен роботов на одном складе~\cite{wurman2008kiva}. В области автономного транспорта группы беспилотных автомобилей требуют согласованного планирования маршрутов для предотвращения столкновений. В поисково-спасательных операциях группы дронов или наземных роботов позволяют параллельно обследовать обширные территории, сокращая время реагирования. В сельском хозяйстве флотилии роботов используются для автоматизации посева, мониторинга и сбора урожая~\cite{lavalle2006planning}.

Основными задачами, возникающими при управлении МРС, являются: определение того, какие задачи выполняет каждый робот, и планирование маршрутов движения роботов таким образом, чтобы они не мешали друг другу. Эти задачи формализуются как Multi-Robot Task Assignment (MRTA) и Multi-Agent Path Finding (MAPF) соответственно.

\subsection{Классификация задач управления мульти-роботизированными системами}

Задачи управления МРС можно разделить на два основных класса, каждый из которых отвечает на свой вопрос~\cite{stern2019mapf, gerkey2004taxonomy}:

\begin{itemize}
    \item \textbf{MRTA} (Multi-Robot Task Assignment) --- задача распределения задач: определяет, \textit{какой} робот выполняет \textit{какую} задачу. Результатом является назначение задач роботам с минимизацией выбранной целевой функции (суммарного расстояния, времени и т.\,д.);
    \item \textbf{MAPF} (Multi-Agent Path Finding) --- задача поиска бесконфликтных маршрутов: определяет, \textit{как} каждый робот добирается до назначенной цели, обеспечивая отсутствие столкновений между агентами. Результатом является набор путей для всех агентов.
\end{itemize}

В реальных системах эти задачи тесно связаны: сначала MRTA-алгоритм назначает задачи роботам, затем MAPF-алгоритм строит бесконфликтные маршруты к назначенным целям. При этом качество решения задачи распределения существенно зависит от навигационных возможностей среды, а эффективность маршрутизации определяется назначенными целевыми позициями.

\subsection{Задача распределения задач в мульти-роботизированных системах}

\subsubsection{Формальная постановка задачи MRTA}

Задача MRTA формулируется следующим образом. Пусть дано множество роботов $R = \{r_1, r_2, \ldots, r_n\}$ и множество задач $T = \{t_1, t_2, \ldots, t_m\}$. Каждая задача $t_j$ характеризуется целевой позицией на карте. Для каждой пары робот--задача определена функция стоимости $c(r_i, t_j)$, отражающая затраты на выполнение задачи $t_j$ роботом $r_i$ (как правило, это расстояние между текущей позицией робота и целевой позицией задачи).

Требуется найти назначение $\sigma$, сопоставляющее каждому роботу подмножество задач, минимизируя суммарную стоимость:

\begin{equation}
    \label{eq:mrta_objective}
    \min_{\sigma} \sum_{i=1}^{n} \sum_{t_j \in \sigma(r_i)} c(r_i, t_j),
\end{equation}

при условии, что каждая задача назначена ровно одному роботу: $\sigma(r_i) \cap \sigma(r_k) = \emptyset$ для $i \neq k$ и $\bigcup_{i=1}^{n} \sigma(r_i) = T$.

\subsubsection{Классификация MRTA-алгоритмов}

Б. П. Джерки и М. Дж. Матарич~\cite{gerkey2004taxonomy} предложили таксономию MRTA-задач по трём измерениям:

\begin{itemize}
    \item \textbf{ST vs MT} (Single-Task vs Multi-Task) --- может ли робот выполнять одну или несколько задач одновременно;
    \item \textbf{SR vs MR} (Single-Robot vs Multi-Robot) --- требуется ли для выполнения задачи один или несколько роботов;
    \item \textbf{IA vs TA} (Instantaneous Assignment vs Time-extended Assignment) --- выполняется ли назначение однократно или с учётом временной динамики.
\end{itemize}

Г. А. Корса, Э. Стенц и М. Б. Диас~\cite{korsah2013taxonomy} расширили данную таксономию, добавив межзадачные зависимости и утилитарные функции. В рамках настоящей работы рассматривается наиболее распространённый класс задач ST-SR --- каждый робот выполняет по одной задаче за раз, и каждая задача требует одного робота.

Помимо этого, различают два режима назначения задач:

\begin{itemize}
    \item \textbf{Online-режим} --- робот получает одну задачу; после её завершения планировщик назначает следующую. Это соответствует варианту IA;
    \item \textbf{Queue-режим} --- планировщик сразу выстраивает полную цепочку задач для каждого робота, что соответствует варианту TA.
\end{itemize}

По подходу к решению MRTA-алгоритмы подразделяются на:

\begin{itemize}
    \item \textit{оптимизационные} --- формулируют задачу как задачу оптимизации и находят точное или приближённое решение (венгерский алгоритм, алгоритм минимальной стоимости потока);
    \item \textit{эвристические} --- используют простые правила для быстрого построения допустимого назначения (жадный алгоритм);
    \item \textit{рыночные} (market-based) --- моделируют процесс торгов, в котором роботы конкурируют за задачи~\cite{dias2006market} (последовательный аукцион, комбинаторный аукцион).
\end{itemize}

\subsubsection{Венгерский алгоритм}

Венгерский алгоритм~\cite{kuhn1955hungarian} решает задачу о назначениях --- частный случай MRTA, в котором число роботов равно числу задач и требуется найти взаимно однозначное соответствие с минимальной суммарной стоимостью.

Формально, задача о назначениях записывается как:

\begin{equation}
    \label{eq:hungarian}
    \min_{\sigma \in S_n} \sum_{i=1}^{n} c_{i, \sigma(i)},
\end{equation}

где $S_n$ --- множество всех перестановок $\{1, 2, \ldots, n\}$, а $c_{i,j}$ --- элемент матрицы стоимостей $C \in \mathbb{R}^{n \times n}$.

Алгоритм основан на теореме Кёнига и последовательно приводит матрицу стоимостей к виду, в котором оптимальное назначение соответствует набору нулевых элементов, покрывающему все строки и столбцы. Вычислительная сложность алгоритма составляет $O(n^3)$, что делает его эффективным для задач умеренного размера.

В контексте МРС венгерский алгоритм работает в online-режиме: за одну итерацию каждому роботу назначается ровно одна задача, а после завершения задачи алгоритм перезапускается для оставшихся незавершённых задач. Матрица стоимостей формируется на основе манхэттенского расстояния между текущими позициями роботов и целевыми позициями задач.

\subsubsection{Жадный алгоритм}

Жадный алгоритм назначения задач представляет собой простейший эвристический подход к решению задачи MRTA. На каждом шаге алгоритм выбирает пару робот--задача $(r_i, t_j)$ с минимальной стоимостью $c(r_i, t_j)$ среди всех ещё не назначенных задач и свободных роботов, и назначает задачу $t_j$ роботу $r_i$.

Вычислительная сложность жадного алгоритма составляет $O(n \cdot m)$, где $n$ --- число роботов, $m$ --- число задач. Алгоритм не гарантирует оптимальности решения, однако обеспечивает быстрое построение допустимого назначения, что может быть критично в системах реального времени.

Аналогично венгерскому алгоритму, жадный алгоритм работает в online-режиме: после завершения задачи роботом алгоритм перезапускается для назначения следующей задачи.

\subsubsection{Алгоритм последовательного аукциона}

Алгоритм последовательного аукциона (Sequential Auction)~\cite{dias2006market} реализует рыночный подход к распределению задач. Задачи выставляются на торги поочерёдно, и роботы подают заявки (ставки) на их выполнение. Задача назначается роботу с минимальной ставкой.

В реализации, рассматриваемой в данной работе, алгоритм работает в queue-режиме и строит полные цепочки задач для каждого робота. Ставка робота рассчитывается как стоимость добавления задачи в конец текущей цепочки: расстояние от виртуальной позиции (положения, в котором робот окажется после выполнения последней задачи в цепочке) до новой задачи. После назначения задачи виртуальная позиция робота обновляется.

Порядок выставления задач определяется разницей между двумя лучшими ставками (bid gap): первыми разыгрываются задачи, по которым расхождение ставок максимально, что уменьшает вероятность субоптимального назначения.

\subsubsection{Алгоритм минимальной стоимости потока}

Алгоритм минимальной стоимости потока (Min-Cost Flow) формулирует задачу MRTA как задачу нахождения потока минимальной стоимости в сети~\cite{ahuja1993network}. Строится двудольный граф, в котором вершины одной доли представляют роботов, а вершины другой --- задачи. Рёбра имеют пропускную способность 1 и стоимость, равную расстоянию между роботом (или его виртуальной позицией) и задачей.

В реализации используется итеративный подход для построения цепочек задач (queue-режим): на каждой итерации решается задача минимальной стоимости потока для текущего <<слоя>> назначений, после чего виртуальные позиции роботов обновляются, и процесс повторяется для оставшихся задач. Решение каждой подзадачи потока выполняется алгоритмом последовательного кратчайшего пути (Successive Shortest Paths).

\subsubsection{Комбинаторный аукцион}

Комбинаторный аукцион (Combinatorial Auction) является расширением последовательного аукциона, в котором роботы подают заявки не на отдельные задачи, а на наборы (связки) задач~\cite{dias2006market}. Это позволяет учитывать синергию между задачами: выполнение двух близко расположенных задач одним роботом может быть более эффективным, чем их назначение разным роботам.

В реализации, рассматриваемой в данной работе, алгоритм работает в два этапа. На первом этапе используется жадное построение цепочек на основе сожаления (regret-based greedy): для каждой задачи вычисляется разница между стоимостью лучшего и второго лучшего назначения, и задачи с максимальным <<сожалением>> назначаются в первую очередь. На втором этапе выполняется локальная оптимизация: алгоритм пытается улучшить решение путём перемещения задач между цепочками роботов и обмена задачами между парами роботов, если это уменьшает суммарную стоимость.

\subsection{Задача поиска бесконфликтных маршрутов}

\subsubsection{Формальная постановка задачи MAPF}

Задача Multi-Agent Path Finding (MAPF)~\cite{stern2019mapf} формулируется следующим образом. Пусть дан неориентированный граф $G = (V, E)$, представляющий среду (в случае сеточной карты вершины соответствуют свободным клеткам, а рёбра --- допустимым переходам между соседними клетками). Задано множество агентов $A = \{a_1, a_2, \ldots, a_k\}$, для каждого из которых определена начальная позиция $s_i \in V$ и целевая позиция $g_i \in V$.

Путь агента $a_i$ представляет собой последовательность вершин $\pi_i = (v_0^i, v_1^i, \ldots, v_{T_i}^i)$, где $v_0^i = s_i$, $v_{T_i}^i = g_i$, и для каждого $t \in \{0, \ldots, T_i - 1\}$ выполняется $(v_t^i, v_{t+1}^i) \in E$ (агент перемещается по ребру) или $v_t^i = v_{t+1}^i$ (агент ожидает на месте). Здесь $T_i$ --- время прибытия агента $a_i$ в целевую вершину.

В задаче MAPF используются два типа конфликтов:

\begin{itemize}
    \item \textbf{Вершинный конфликт} (vertex conflict) --- два агента $a_i$ и $a_j$ занимают одну вершину в один момент времени: $v_t^i = v_t^j$ при $i \neq j$;
    \item \textbf{Рёберный конфликт} (edge conflict, swap conflict) --- два агента обмениваются позициями: $v_t^i = v_{t+1}^j$ и $v_{t+1}^i = v_t^j$ при $i \neq j$.
\end{itemize}

Набор путей $\Pi = \{\pi_1, \ldots, \pi_k\}$ называется \textit{решением} задачи MAPF, если он не содержит конфликтов. Для оценки качества решения используются две основные метрики:

\begin{equation}
    \label{eq:makespan}
    \text{Makespan} = \max_{i \in \{1, \ldots, k\}} T_i,
\end{equation}

\begin{equation}
    \label{eq:soc}
    \text{SOC (Sum of Costs)} = \sum_{i=1}^{k} T_i.
\end{equation}

Makespan отражает общее время до завершения миссии всеми агентами, тогда как SOC отражает суммарные затраты ресурсов (шагов) всех агентов. Задача нахождения оптимального решения MAPF (по любой из метрик) является NP-трудной~\cite{yu2013mapf_complexity}, что мотивирует разработку как точных алгоритмов с приемлемой практической сложностью, так и приближённых алгоритмов.

\subsubsection{Классификация MAPF-алгоритмов}

Существующие MAPF-алгоритмы классифицируются по нескольким критериям~\cite{stern2019mapf}:

\begin{enumerate}[arabic]
    \item \textbf{По архитектуре управления:}
    \begin{itemize}
        \item \textit{централизованные} --- единый планировщик строит пути для всех агентов одновременно, располагая полной информацией о среде (CBS, ECBS, EECBS, M*);
        \item \textit{децентрализованные} --- агенты планируют пути независимо или последовательно, обмениваясь ограниченной информацией (CA*, WHCA*, PIBT).
    \end{itemize}
    \item \textbf{По оптимальности:}
    \begin{itemize}
        \item \textit{оптимальные} --- гарантируют нахождение решения с минимальной стоимостью (CBS, M*);
        \item \textit{ограниченно субоптимальные} --- гарантируют, что стоимость решения не превышает оптимальную более чем в $w$ раз (ECBS, EECBS);
        \item \textit{неоптимальные} --- не дают гарантий оптимальности, но работают значительно быстрее (CA*, WHCA*, PIBT).
    \end{itemize}
    \item \textbf{По режиму работы:}
    \begin{itemize}
        \item \textit{offline} --- все пути строятся до начала движения агентов (CBS, ECBS, EECBS, M*, CA*, WHCA*);
        \item \textit{online} --- пути строятся или корректируются в процессе движения агентов (PIBT).
    \end{itemize}
\end{enumerate}

В таблице~\ref{tab:mapf_classification} приведена сводная классификация реализованных алгоритмов.

\begin{table}
    \caption{Классификация реализованных MAPF-алгоритмов}
    \begin{tabularx}{\textwidth}{|X|c|c|c|}\hline
    Алгоритм & Архитектура & Оптимальность & Режим \\ \hline
    CBS (A*)     & централиз. & оптимальный         & offline \\ \hline
    CBS (SIPP)   & централиз. & оптимальный         & offline \\ \hline
    ECBS         & централиз. & огр. субоптим.      & offline \\ \hline
    EECBS        & централиз. & огр. субоптим.      & offline \\ \hline
    M*           & централиз. & оптимальный         & offline \\ \hline
    CA*          & децентрал. & неоптимальный       & offline \\ \hline
    WHCA*        & децентрал. & неоптимальный       & offline \\ \hline
    PIBT         & децентрал. & неоптимальный       & online  \\ \hline
    \end{tabularx}
    \label{tab:mapf_classification}
\end{table}

\subsubsection{Алгоритм A* с расширением по времени}

Алгоритм A*~\cite{hart1968astar} является фундаментальным алгоритмом поиска кратчайшего пути в графе и лежит в основе большинства MAPF-алгоритмов. В контексте задачи MAPF используется его модификация с расширением по времени (time-expanded A*), в которой состояние агента определяется не только позицией, но и моментом времени: $(v, t)$, где $v \in V$ --- текущая вершина, $t \in \mathbb{N}_0$ --- текущий временной шаг.

Функция оценки каждого состояния определяется как:

\begin{equation}
    \label{eq:astar}
    f(v, t) = g(v, t) + h(v),
\end{equation}

где $g(v, t) = t$ --- стоимость достижения состояния $(v, t)$ от начального состояния $(s_i, 0)$,\par $h(v)$ --- эвристическая оценка стоимости пути от $v$ до цели $g_i$.

В качестве эвристики используется манхэттенское расстояние:

\begin{equation}
    \label{eq:manhattan}
    h(v) = |x_v - x_{g_i}| + |y_v - y_{g_i}|,
\end{equation}

где $(x_v, y_v)$ и $(x_{g_i}, y_{g_i})$ --- координаты вершины $v$ и целевой вершины $g_i$ соответственно.

Алгоритм A* с расширением по времени поддерживает множество ограничений (reservations) --- набор состояний $(v, t)$, которые агент не может занимать (зарезервированы другими агентами). Это позволяет планировать пути с учётом уже построенных путей других агентов.

\subsubsection{Conflict-Based Search (CBS)}

Алгоритм CBS (Conflict-Based Search)~\cite{sharon2015cbs} является оптимальным централизованным алгоритмом решения задачи MAPF. CBS использует двухуровневую архитектуру поиска.

\textbf{Верхний уровень} оперирует деревом ограничений (Constraint Tree, CT). Каждый узел CT содержит:
\begin{itemize}
    \item набор ограничений $\mathcal{C} = \{C_1, C_2, \ldots, C_k\}$, где $C_i$ --- множество ограничений для агента $a_i$;
    \item набор путей $\Pi = \{\pi_1, \pi_2, \ldots, \pi_k\}$, согласованных с ограничениями;
    \item суммарную стоимость $\text{cost} = \sum_{i=1}^{k} |\pi_i|$.
\end{itemize}

Корень дерева содержит пустые множества ограничений. На каждом шаге верхний уровень извлекает узел CT с минимальной стоимостью и проверяет наличие конфликтов в наборе путей. Если конфликтов нет, решение найдено.

При обнаружении конфликта между агентами $a_i$ и $a_j$ (вершинного: $(a_i, a_j, v, t)$, или рёберного: $(a_i, a_j, u, v, t)$) узел разветвляется на два дочерних узла. В первом добавляется ограничение для агента $a_i$ (запрет на нахождение в вершине $v$ в момент $t$ или на переход $(u, v)$ в момент $t$), во втором --- аналогичное ограничение для агента $a_j$.

\textbf{Нижний уровень} решает задачу поиска кратчайшего пути для одного агента с учётом заданного набора ограничений. В качестве нижнеуровневого планировщика используется алгоритм A* с расширением по времени, описанный выше.

CBS гарантирует оптимальность решения (по метрике SOC) и полноту. Вычислительная сложность в худшем случае экспоненциальна относительно числа агентов, однако на практике алгоритм эффективен для сценариев с умеренным числом конфликтов.

\textbf{Вариант CBS с SIPP.} В качестве альтернативного нижнеуровневого планировщика может использоваться алгоритм SIPP (Safe Interval Path Planning)~\cite{phillips2011sipp}. SIPP оперирует безопасными временными интервалами --- промежутками времени, в течение которых данная вершина свободна от ограничений. Состояние поиска определяется как $(v, I)$, где $I$ --- безопасный интервал для вершины $v$. Это позволяет существенно сократить пространство поиска по сравнению с time-expanded A* при большом количестве ограничений, так как каждый интервал агрегирует множество временных шагов.

\subsubsection{Enhanced CBS (ECBS)}

Алгоритм ECBS (Enhanced CBS)~\cite{barer2014ecbs} является ограниченно субоптимальной модификацией CBS, обеспечивающей значительное ускорение за счёт ослабления требования оптимальности.

ECBS вводит параметр субоптимальности $w \geq 1$ и гарантирует, что стоимость найденного решения не превышает оптимальную более чем в $w$ раз:

\begin{equation}
    \label{eq:ecbs_bound}
    \text{cost}(\Pi_{\text{ECBS}}) \leq w \cdot \text{cost}(\Pi^*).
\end{equation}

Ускорение достигается за счёт использования фокального поиска (focal search) на обоих уровнях алгоритма. На верхнем уровне поддерживается два списка открытых узлов: основной (OPEN), упорядоченный по стоимости, и фокальный (FOCAL), содержащий узлы со стоимостью не более $w \cdot f_{\min}$, где $f_{\min}$ --- минимальная стоимость в OPEN. Узлы из FOCAL выбираются по вторичному критерию, например, по числу конфликтов.

На нижнем уровне аналогично используется фокальный A*: среди путей, стоимость которых не превышает $w \cdot f_{\min}^{\text{low}}$, выбирается путь с минимальным числом конфликтов с уже построенными путями других агентов. Это приводит к построению путей, порождающих меньше конфликтов на верхнем уровне, что сокращает размер дерева ограничений.

\subsubsection{Explicit Estimation CBS (EECBS)}

Алгоритм EECBS (Explicit Estimation CBS)~\cite{li2021eecbs} является дальнейшим развитием ECBS, повышающим эффективность фокального поиска на верхнем уровне за счёт использования онлайн-эвристики.

В ECBS выбор узлов из фокального списка основывается на числе конфликтов, которое лишь косвенно связано с итоговой стоимостью решения. EECBS вводит явную оценку стоимости разрешения конфликтов: для каждого узла CT вычисляется $\hat{h}$ --- оценка дополнительной стоимости, необходимой для устранения всех оставшихся конфликтов. Эта оценка обновляется в процессе поиска на основе информации, полученной при разветвлении узлов.

Узлы фокального списка упорядочиваются по $\hat{f} = g + \hat{h}$, что направляет поиск в сторону узлов, более вероятно приводящих к решению с низкой стоимостью. Граница субоптимальности $w$ сохраняется, однако на практике EECBS находит решения значительно ближе к оптимуму, чем ECBS с тем же значением $w$.

\subsubsection{M* (Subdimensional Expansion)}

Алгоритм M*~\cite{wagner2015mstar} решает задачу MAPF путём поиска в совместном пространстве состояний, но в отличие от наивного подхода выполняет расширение в полное пространство состояний только при возникновении конфликтов.

В совместном пространстве состояний конфигурация системы из $k$ агентов представляется как кортеж $(v_1, v_2, \ldots, v_k) \in V^k$, а переход --- как одновременное перемещение всех агентов. Наивный поиск в таком пространстве имеет сложность $O(|V|^k)$, что делает его непрактичным для большого числа агентов.

M* использует принцип поддимерного расширения (subdimensional expansion): каждому агенту предварительно вычисляется индивидуальная оптимальная политика (кратчайший путь без учёта других агентов) с помощью обратного алгоритма Дейкстры от целевой вершины. При поиске в совместном пространстве каждый агент по умолчанию следует своей индивидуальной политике. Когда обнаруживается конфликт между агентами $a_i$ и $a_j$ в некоторой конфигурации, формируется множество столкновений (collision set), и конфигурация расширяется по всем возможным комбинациям действий конфликтующих агентов.

Множества столкновений обратно распространяются (back-propagated) к предшествующим конфигурациям, чтобы обеспечить полноту поиска. Алгоритм M* является оптимальным и полным. Его эффективность определяется степенью связности конфликтов: если конфликты затрагивают лишь небольшие подмножества агентов, M* работает значительно быстрее, чем поиск в полном совместном пространстве.

\subsubsection{Cooperative A* (CA*)}

Алгоритм Cooperative A* (CA*)~\cite{silver2005cooperative} реализует децентрализованный подход к решению задачи MAPF. Агенты планируют пути последовательно, один за другим, с использованием общей таблицы резерваций.

Алгоритм работает следующим образом:
\begin{enumerate}[arabic]
    \item агенты упорядочиваются в некотором порядке (например, по приоритету);
    \item для первого агента строится кратчайший путь с помощью A* с расширением по времени; его путь добавляется в таблицу резерваций;
    \item для каждого следующего агента строится путь с учётом всех ранее зарезервированных позиций;
    \item целевая позиция агента, достигшего цели, также резервируется на все последующие временные шаги.
\end{enumerate}

Таблица резерваций содержит множество троек $(v, t)$, запрещающих нахождение очередного агента в вершине $v$ в момент $t$. Помимо вершинных резерваций, учитываются и рёберные --- запрет на переход по ребру $(u, v)$ в момент $t$, если другой агент проходит по ребру $(v, u)$ в тот же момент.

CA* не гарантирует оптимальности решения, поскольку результат зависит от порядка планирования агентов: агенты, планирующие первыми, имеют больше свободы выбора пути. Более того, в некоторых конфигурациях CA* может не найти решение, даже если оно существует, если порядок агентов приводит к тупиковой ситуации. Тем не менее, CA* обладает низкой вычислительной сложностью --- $O(k \cdot |V| \cdot T)$, где $T$ --- максимальная глубина по времени, --- и хорошо масштабируется при большом числе агентов.

\subsubsection{Windowed Hierarchical CA* (WHCA*)}

Алгоритм WHCA*~\cite{silver2005cooperative} является модификацией CA*, в которой планирование ограничивается фиксированным временным окном $W$.

На каждой итерации WHCA* каждый агент планирует путь на ближайшие $W$ временных шагов с использованием алгоритма CA* (последовательное планирование с таблицей резерваций). По истечении окна все агенты перепланируют пути от текущих позиций на следующие $W$ шагов.

Ограничение горизонта планирования имеет два ключевых преимущества:
\begin{itemize}
    \item уменьшается пространство поиска на каждой итерации, что ускоряет планирование;
    \item агенты, ещё не находящиеся вблизи друг друга, не создают избыточных ограничений.
\end{itemize}

Недостатком WHCA* является то, что локально оптимальные решения на каждом окне могут привести к глобально неоптимальным или даже тупиковым ситуациям. Размер окна $W$ является параметром, определяющим баланс между качеством решения и скоростью планирования.

\subsubsection{Priority Inheritance with Backtracking (PIBT)}

Алгоритм PIBT~\cite{okumura2022pibt} представляет собой онлайн-децентрализованный алгоритм, в котором агенты планируют перемещение всего на один временной шаг вперёд. Это делает PIBT самым <<реактивным>> из рассматриваемых алгоритмов.

Каждому агенту назначается приоритет, который определяет порядок планирования на текущем шаге. Приоритеты обновляются динамически: агент, не достигший цели, получает более высокий приоритет, чем агент, уже находящийся в целевой вершине.

Ключевым механизмом алгоритма является наследование приоритетов (priority inheritance): если агент $a_i$ с высоким приоритетом хочет занять вершину, в которой находится агент $a_j$ с более низким приоритетом, то $a_j$ <<наследует>> приоритет $a_i$ и вынужден переместиться. Если $a_j$ не может найти свободную соседнюю вершину, инициируется откат (backtracking): $a_i$ отменяет свой запрос и пробует другое направление.

Для выбора предпочтительного направления движения каждый агент использует таблицу расстояний (distance table) --- матрицу кратчайших расстояний от каждой вершины до целевой, вычисленную предварительно обратным поиском в ширину.

PIBT не гарантирует оптимальности решения и полноты в общем случае, однако обеспечивает быстрое реактивное управление и хорошо масштабируется при большом числе агентов. Вычислительная сложность одного шага составляет $O(k)$, где $k$ --- число агентов.

\subsection{Интеграция MRTA и MAPF}

В реальных мульти-роботизированных системах задачи распределения и навигации решаются совместно: MRTA-алгоритм определяет, какие задачи выполняет каждый робот, а MAPF-алгоритм строит бесконфликтные маршруты к назначенным целям. Этот двухэтапный подход представляет собой основную архитектурную модель, реализованную в рассматриваемом программном комплексе.

При online-режиме назначения каждый робот получает одну задачу, после чего все активные роботы совместно планируют маршруты через MAPF-алгоритм. По завершении задачи роботом планировщик назначает ему следующую задачу, и маршрут для данного робота перестраивается с учётом оставшихся путей других агентов.

При queue-режиме каждому роботу сразу назначается полная цепочка задач. Начальные маршруты к первым задачам планируются совместно через выбранный MAPF-алгоритм. По мере завершения задач роботы переходят к следующим элементам цепочки, и маршруты достраиваются последовательно в стиле CA* --- с учётом путей остальных агентов.

Объективное сравнение различных комбинаций MRTA- и MAPF-алгоритмов при варьировании параметров среды (размер карты, число роботов, плотность препятствий, стратегия размещения) является основной задачей разрабатываемого программного комплекса.
